\chapter{Introduction and Background}
\label{cha:intro}

\section{Introduction}
\label{sec:introduction}

The rapid growth of intelligent sensing devices is shifting AI from power-hungry
data centers and supercomputing systems to low-power edge platforms, where
perceptual AI, text generation, reasoning, and control occur near the data source
under challenging latency, energy, and privacy constraints~\cite{reddi2025generative}.

Within the domain of perceptual AI, TinyML was the initial step toward edge
intelligence. It has been demonstrated that inference workloads are feasible on
microcontrollers (MCUs) using moderately parameterized CNN and RNN
models~\cite{10177729,mi13060851}. Since 2024, the edge has evolved to embrace
Generative Edge AI. Modern Edge AI is built on heterogeneously quantized
Transformer architectures, including small language models (SLMs), which seek to
retain useful generative and reasoning capabilities under tight resource (memory
and computation) budgets~\cite{technologies12060081,electronics14173468}.
Techniques such as deep quantization, pruning, knowledge distillation, and
hardware-aware neural architecture design have become central to reconciling model
accuracy with edge deployability~\cite{s25103191}. Delivering real-time
responsiveness requires tight integration of models, algorithms, and low-power
hardware.

In this landscape, conversational audio is a natural target for edge deployment:
its pipeline decomposes into independently optimizable modular stages, its temporal
features operate at comparatively low dimensionality, and its biometric content
makes local processing preferable from both privacy and regulatory perspectives.
Conversational systems impose stringent real-time constraints that benefit from the
elimination of network latency, and their modular structure allows individual
stages to be aggressively quantized without prohibitive quality loss.

This thesis first defines the objectives of the proposed work on conversational
systems for edge devices in Section~\ref{sec:research_question}. Section~\ref{sec:sota}
provides an overview of the state of the art in modern AI architectures, edge AI
systems, and embedded speech processing. Section~\ref{sec:contribution} outlines the
specific findings and contributions of this work. The methodology for achieving
conversational-quality inference under strict compute and memory budgets is detailed
in Chapter~\ref{cha:design}, while Chapter~\ref{cha:experiments} presents the
experimental setup, the evaluation, and the discussion of the broader implications
and limitations of the findings. Chapter~\ref{cha:conclusioni} concludes the work
and outlines future directions.

\section{Research Question}
\label{sec:research_question}

Conversational AI systems aim to provide rich, context-aware interactions through
natural spoken dialogue. Achieving this capability requires models that can process
audio input, understand and reason upon user intentions, and generate appropriate
responses within strict real-time constraints. At the same time, conversational
audio often contains sensitive and biometric information, making local processing
desirable from both privacy and reliability perspectives.

Conversational speech interfaces operate within defined latency
thresholds~\cite{doherty2015flow,defossez2025moshi}: response times of
100--200\,ms are perceived as instantaneous, delays of 0.5--1\,s are categorized as
immediate and remain acceptable for sustained dialogue, and intervals exceeding
1\,s degrade the perceived interaction quality. Natural human turn-taking averages
approximately 230\,ms~\cite{defossez2025moshi}, placing the acceptable
machine-response boundary near the transition between the instantaneous and
immediate bands. This leaves a narrow window for the entire pipeline (audio
encoding, ASR, language modeling, and text-to-speech synthesis) to complete on a
per-utterance basis.

Edge devices operate under strict limits in compute capability, memory footprint,
and energy consumption, typically drawing between 3 and 10\,W of power for embedded
platforms in the smart-assistant and IoT category. Yet meaningful conversational
interaction relies on resource-intensive transformer-based language models whose
autoregressive decoding generates one token at a time, with each token requiring
tens to hundreds of billions of multiply-accumulate operations~\cite{pope2022efficiently}.

A viable edge deployment must therefore satisfy three constraints simultaneously:
\begin{enumerate}[label=(\roman*)]
    \item per-token generation latency within the 200--500\,ms window such that the aggregated pipeline remains under the end-to-end budget,
    \item total system power in the single-digit-watt range compatible with embedded and battery-backed operation,
    \item local-only execution that preserves the privacy of conversational data without reliance on cloud offload.
\end{enumerate}

The central question this work poses is whether a combination of model quantization
and selective hardware acceleration can bring a competitive conversational model
within these constraints on a resource-limited edge platform.

\section{State of the Art}
\label{sec:sota}

\subsection{Evolution of Modern AI Models}
Recent years have seen rapid progress in AI driven by hyper-parameterized neural
topologies trained on tens of trillion-token datasets. Transformer-based models
have become the dominant paradigm across multiple modalities~\cite{10433480}. In
vision, transformer architectures and diffusion models have enabled high-capacity
systems to achieve image understanding and generation~\cite{10902142}. Large
language models (LLMs) have demonstrated strong reasoning, chain-of-thought (CoT)
capabilities, dialogue, and code generation~\cite{10.1145-3747588}. Similarly,
research on speech has increasingly adopted end-to-end neural architectures that
unify acoustic modeling, language modeling, and decoding within a single
framework~\cite{app142411583}. As these models become more capable, they are
increasingly integrated into systems that interact with users and the physical
world in real time.

\subsection{Real-time and Edge AI Systems}
Recent research has focused on integrating AI models into systems that respond to
continuous streams of sensory input. Vision--Language--Action (VLA) architectures
exemplify this trend by integrating visual-audio-tactile perception with
language-based reasoning and control
policies~\cite{huang2026tactilevla,huang2026tafvlatactileforcealignmentvisionlanguageaction,bi2025vlatouchenhancingvisionlanguageactionmodels}.
Applications such as conversational assistants, robotics, and augmented interfaces
require real-time inference with strict latency constraints and reliable operation
in time-varying environments. These requirements motivated increased interest in
edge-based inference, where processing occurs close to the sensing element,
eliminating cloud delay concerns. Several methods have therefore been proposed to
deploy machine learning models on the edge, operating under limited compute, memory,
and energy budgets. From the early work on TinyML, more recent approaches harness
model compression techniques such as quantization, pruning, and knowledge
distillation, as well as hardware-aware neural architecture
design~\cite{dnn_compression}. Together, these advances form the foundation of
modern Edge AI, which seeks to balance model capability with the strict resource
constraints of edge platforms.

\subsection{Speech and Conversational Architectures}
De-facto standard speech processing follows a modular pipeline in which an ASR
module converts audio to text, a natural language understanding component extracts
intent, a language model generates a textual response, and a TTS engine synthesizes
the output waveform. This decomposition enables each stage to be optimized
independently for latency, accuracy, and power, and is the dominant architecture in
deployed conversational systems. Typical ASR models range from compact encoder-only
architectures (Whisper-tiny at 39\,M parameters) to large encoder-decoder designs
(Whisper-large-v3 at 1.55\,B parameters)~\cite{radford2023whisper}, while TTS
engines span 10\,M to 500\,M parameters. Streaming ASR models allow speech to be
transcribed incrementally as audio is received, reducing the time to first token
(TTFT) compared to batch-based processing. The aggregate parameter count across
three or four separate models can reach several billion, and the per-stage inference
latencies sum within the end-to-end response budget of 200--500\,ms.

Recent work has explored end-to-end architectures that unify these stages into a
single neural model. LFM2-Audio~\cite{amini2025lfm2technicalreport} represents this
trend well: a 1.5\,B-parameter model that jointly handles audio encoding, language
modeling, and speech synthesis within one framework, eliminating inter-stage latency
and easing the deployment footprint. LFM2.5 denotes a subsequent release of the same
architecture with updated weights; the 1.2\,B language backbone used in this work is
architecturally identical across both releases. The deployment requirements of such
models vary substantially; a quantitative profiling of representative checkpoints,
presented in Section~\ref{subsec:profiling}, motivates the specific model choice
adopted in this work. While end-to-end architectures simplify the system design,
their computational cost remains concentrated in a single large model rather than
distributed across independently scalable stages, making hardware acceleration of the
dominant internal kernel an attractive optimization target.

\subsection{Edge Speech Systems and On-Device Conversational AI}
The characteristics of speech signals make audio workloads well suited for edge
deployment. Lightweight components such as wake-word detection and voice activity
detection can run continuously on MCUs burning power under 1\,W, while more demanding
pipeline stages (ASR, language modeling) are typically hosted on single-board
computers or heterogeneous SoCs in the 3--15\,W range. The Raspberry Pi 4B
(quad-core Cortex-A72 at 1.5\,GHz, up to 8\,GB RAM) has emerged as a common
evaluation platform for on-device speech systems, hosting engines such as
Faster-Whisper for real-time voice interaction~\cite{10.1145/3776942.3776981}. At the
lower end of the resource spectrum, comparative studies of phoneme-based and
word-based recognition across MCU-class devices (Cortex-M4/M7 at 64--480\,MHz, with
hundreds of kilobytes of SRAM) identify specific architectures that optimize the
accuracy-versus-memory trade-off for offline keyword spotting~\cite{10711025}.
Prototype systems integrating Vosk or Whisper-based ASR with local NLP microservices
have demonstrated sub-second speech-to-text latency entirely on edge
hardware~\cite{info16080685}, validating the feasibility of local-only processing for
conversational audio.

Despite this progress, existing works have largely addressed individual pipeline
components in isolation. Comparatively little research examines conversational
interaction as an end-to-end systems problem in which the language backbone, the ASR
front-end, and the TTS output stage must all complete within the latency and power
budget defined in Section~\ref{sec:research_question}. Addressing this challenge
requires not only lightweight models but also system-level optimization spanning
streaming inference, hardware-software partitioning, and selective acceleration of
the dominant computational kernels.

\subsection{Profiling of Speech Models for Edge Deployment}
\label{subsec:profiling}
To evaluate the feasibility of conversational inference on edge devices, a
representative set of speech and multimodal models was profiled using
gguf-parser-go\footnote{\texttt{https://github.com/gpustack/gguf-parser-go}}, a tool
that extracts structural and resource information from GGUF-format checkpoints. The
profiling covers parameter count, memory footprint, memory bandwidth requirements,
and estimated compute throughput. These are the same metrics used in a recent
large-scale study of 21,039 checkpoints across six workload
categories~\cite{sym18050766}, which found that audio-language models define the
lightest deployment envelope (mean 13.41\,GB/s bandwidth, 134\,GFLOPS compute,
826\,MiB UMA). One observation was that parameter count alone is an insufficient
predictor of edge feasibility. The resulting comparison, illustrated in
Fig.~\ref{fig:efficiency}, reveals wide variation in deployment requirements:
parameter counts span from 500\,M (OuteTTS-0.2-500M-Q4) to 9.4\,B
(THUDM\_GLM-4-9B-0414-Q4), RAM footprints range from 0.29\,GB (LFM2.5-Audio-1.5B-Q4,
the lowest among all profiled models) to 2.67\,GB (Qwen2-7B-LLM-Q4), and estimated
memory bandwidth demand varies from 3.6\,GB/s to 65\,GB/s. Several architectures,
including the 7--9\,B parameter class, exceed practical memory and bandwidth limits
for embedded platforms despite offering strong conversational capabilities. In
contrast,
LFM2.5-Audio-1.5B\footnote{\texttt{https://huggingface.co/LiquidAI/LFM2.5-Audio-1.5B}}
combines the lowest RAM footprint with moderate bandwidth requirements (3.9\,GB/s) and
a competitive parameter scale (1.18\,B), positioning it closest to the desired
resource envelope of embedded conversational devices.

\begin{figure}[!ht]
  \centering
  \includegraphics[width=\linewidth]{audio_benchmark_valori_reali.png}
  \caption{Model throughput (TOPS) as a function of resource requirements. Models with the highest throughput at the lowest resource cost appear toward the bottom left.}
  \label{fig:efficiency}
\end{figure}

Based on this analysis, LFM2.5 was chosen as the primary candidate for the subsequent
hardware and software studies, as it lies closest to the efficiency--capability
trade-off required for conversational inference under edge resource constraints.

\subsection{Related Work on the Kria KV260 Platform}
Following the identification of LFM2.5 as a suitable candidate for edge deployment,
the next step was to map its computational workload onto an appropriate hardware
platform. Field-Programmable Gate Arrays (FPGAs) provide a compelling and flexible
target for edge AI acceleration due to their ability to exploit fine-grained
parallelism, customize dataflow, and operate within tight power and latency budgets.

Recent work on the AMD-Xilinx Kria KV260 demonstrated the feasibility of
accelerating transformer models. The three designs surveyed below target both
feed-forward network (FFN) and self-attention matrix multiplications, which together
accounted for the greater part of the inference latency and bandwidth consumption in
transformer workloads. A 6.6k-parameter Tiny Transformer for ECG arrhythmia
detection deployed on the KV260 achieved 56.62\,ms total inference latency at
77\,MHz, with the FFN contributing 54.9\% of runtime (31.11\,ms)~\cite{11310944tin}.
A tiled matrix multiplication accelerator targeting the Q, K, and V self-attention
projections of DistilBERT achieved a $7\times$ speedup over the Arm CPU baseline at
100\,MHz and 3.3\,W, delivering 3.12\,GFLOPS of standalone GEMM throughput on the
KV260~\cite{li2025fpgabasedhardware}. A BRAM-banked double-buffered matrix
multiplication accelerator on the same platform reached a peak compute throughput of
22.83\,GOPS for a DistilBERT-scale FFN multiplication ($7.31\times$ over its
baseline), completing the operation in 13.2\,ms at 100\,MHz. The study identified
PS-side data handling as the dominant system-level bottleneck, a finding consistent
with the orchestration overheads observed in the present work
(Section~\ref{sec:discussion})~\cite{11311069}.

Existing approaches generally focus on accelerating full transformer blocks or
large-scale models, where attention and feed-forward computations are treated
together as co-dominant workloads. This work targets the minimal set of computational
kernels required to achieve conversational-level inference, focusing on feed-forward
acceleration. The profiling results presented in Section~\ref{subsec:cpuexec} show
that, for small models such as LFM2.5-1.2B, the execution cost is overwhelmingly
dominated by feed-forward components, with attention contributing a comparatively
smaller fraction of the total runtime. This shift in computational balance motivated a
different design perspective, in which the feed-forward path becomes the primary
target for acceleration. Furthermore, the LFM2.5-1.2B variant operates using quantized
weights without bias terms, simplifying the arithmetic structure of the model and
making it particularly well suited for efficient FPGA implementation.

\section{Contributions}
\label{sec:contribution}
This work provides five contributions.

\begin{enumerate}
    \item \textbf{System-level bottleneck identification for edge conversational inference.}
    Fine-grained profiling of LFM2.5-1.2B on an Arm-based edge platform shows that
    feed-forward MatMul + SwiGLU computation dominates runtime, motivating selective
    acceleration instead of full-model hardware mapping.

    \item \textbf{A practical CPU--FPGA offload integration in the llama.cpp/ggml stack.}
    A fused FFN hardware operator is implemented with graph-level interception,
    layer-aware weight caching, interrupt-driven synchronization, and fallback
    compatibility with standard CPU execution.

    \item \textbf{A resource-aware FPGA design and characterization.}
    The accelerator is implemented as a staged dataflow pipeline and analyzed at
    stage granularity, highlighting where execution is limited by memory transfer
    across all projection stages, including the down-projection. The design operates
    near device limits (e.g., high LUT occupancy), exposing realistic scaling
    constraints on KV260-class platforms.

    \item \textbf{Quantitative efficiency/performance trade-off across thread regimes.}
    CPU Only versus FPGA+CPU measurements across thread counts (T1--T4) are reported
    using throughput, latency, power, and performance per watt. Results show clear
    efficiency benefits at low and moderate CPU thread counts, while high-thread
    regimes remain constrained by data movement and orchestration overheads.

    \item \textbf{Comparative analysis against an edge MPU.}
    The STM32MP2, a low-power dual-core Arm A35 running at 1.5\,GHz, is benchmarked
    against the FPGA+CPU implementation in terms of latency and energy consumption.
\end{enumerate}
